\section{Joint forecasts for observed operational cohorts}
At forecast origin $t$, let $V_t$ be the set of currently active entities and
$\mathcal H_t$ the event history observable by $t$. Each entity has a current
state $s_i(t)$ and zero or more observed relation memberships $g_r(i,t)$.
Memberships may be shared across states. A missing membership creates no shared
factor. Identities index the observed incidence structure; they are not learned
identity embeddings.

For ordered horizons $h_1<\cdots<h_K$, define
$Y_{ik}=\mathbb{1}\{T_i\leq t+h_k\}$, where $T_i$ is the first exit from the
entity's current task state or, in the workflow application, its first qualifying
application-resolution event. We estimate the conditional joint law of these
outcomes for the entire current cohort. This is a discrete-horizon forecast:
each draw assigns an exit bin or survival beyond $h_K$. It does not specify
within-bin event times, subsequent returns, or outcomes for future arrivals.

An individual forecaster supplies a monotone CDF
$F_{ik}=\Pr(T_i\leq t+h_k\mid\mathcal H_t)$.
Individual predictive accuracy does not determine collective uncertainty.
For a known group $A\subseteq V_t$, its exit count is
$C_{Ak}=\sum_{i\in A}Y_{ik}$ and
\begin{equation}
 \operatorname{Var}(C_{Ak}\mid\mathcal H_t)
 =\sum_{i\in A}F_{ik}(1-F_{ik})
 +2\sum_{i<j,\,i,j\in A}\operatorname{Cov}(Y_{ik},Y_{jk}\mid\mathcal H_t).
\end{equation}
This standard decomposition motivates a controlled comparison: keep every
$F_{ik}$ fixed and vary only the joint sampling law.

\subsection{Observed-incidence factor sampling}
For each forecast draw, independently sample a global factor $Z_0$, one factor
$Z_{r,g}$ per observed relation group, and a private factor $\epsilon_i$ per
entity. The Gaussian construction is
\begin{equation}
 Z_i=a_iZ_0+\sum_{r=1}^{R}b_{ir}Z_{r,g_r(i,t)}+c_i\epsilon_i,
 \qquad a_i^2+\sum_r b_{ir}^2+c_i^2=1,
 \qquad U_i=\Phi(Z_i).
\end{equation}
All underlying factors are standard normal. The nonnegative loading
parameterization models positive residual association. It cannot express
general negative dependence. When a relation is unavailable, its squared
loading is transferred to the private component. Consequently, each $U_i$ is
uniform, independently of the cohort size and incidence pattern. Setting
$Y_{ik}^{(m)}=\mathbb{1}\{U_i^{(m)}\leq F_{ik}\}$ preserves the supplied
marginals and makes every sampled trajectory monotone across horizons.
The same shared draws are reused across streamed entity chunks.

We also evaluate a classical common-shock construction related to
Marshall--Olkin models \cite{marshall1967}. With independent unit-rate
exponential factors $E_0,E_{r,g},E_i$ and nonnegative weights summing to one,
\begin{equation}
 L_i=\min\left\{\frac{E_0}{a_i^2},
       \left\{\frac{E_{r,g_r(i,t)}}{b_{ir}^2}\right\}_{r=1}^{R},
       \frac{E_i}{c_i^2}\right\},
 \qquad U_i=1-\exp(-L_i).
\end{equation}
A zero-weight factor has an infinite race time. Since the minimum has unit
exponential rate, this construction also preserves each individual CDF. Neither
construction is a new copula family or an identified causal shock model.

The independent comparison sets the private weight to one. The global model
learns state-specific global/private weights. The static relation model adds
state-specific weights for each relation type. The context model adds a
two-layer, width-32 MLP offset to the state-specific logits, with bounded offsets
and a quadratic shrinkage penalty. A softmax allocates variance weights, with
at least $0.02$ private mass. This prevents perfect dependence in fitted models.
Sampling costs are linear in the number of entities and observed memberships
per draw; no dense entity-by-entity covariance matrix is formed.

\subsection{Fitting without in-sample forecast residuals}
We fit dependence on chronological forward-fold predictions: every training
CDF was produced by an individual model trained before that target period,
with any calibrator fitted in a separate earlier window. At each origin we
sample uniform entity pairs, uniform eligible relation-group pairs, and
group-size-weighted pairs. Sampling uses observed membership and group size,
not future outcomes. Pairs may span current states.

Training minimizes the composite negative log probability of the observed
binary pair, with independently sampled horizon indices for its two entities.
The Gaussian pair law uses the correlation implied by shared factors and a
bivariate normal CDF; the common-shock pair law has an analytic expression.
Validation uses a fixed balanced pattern over the $K^2$ horizon pairs. Pair
likelihood is a tractable fitting objective; it does not guarantee optimal
aggregate scores or operational decisions. Dependence parameters can absorb
individual-model misspecification as well as residual association.

\subsection{A control on relation incidence}
For a fitted static model, we permute relation memberships within current state
and half-unit logit-risk bins of the last-horizon marginal prediction. Known
group sizes within each stratum, all individual CDFs, and fitted loading weights
are preserved. Evaluation groups retain their true observed memberships.
Thus, only the identity of entities sharing relation noise changes. A benefit
from true incidence supports the predictive usefulness of this grouping under
the fitted model; it does not establish that changing an actual job or handler
assignment would cause the same effect.

\section{Evaluation quantities}
For each origin, we score current-state population counts and observed relation
groups containing at least five entities. Scores are normalized by the group's
known current size. We average equally across groups within a projection,
horizons, projection types, and forecast origins. We also report projections
separately so a population-level benefit cannot conceal a relation-level loss.

For $M$ independent scenario draws, the fair ensemble CRPS
\cite{ferro2014,gneiting2007} is
\begin{equation}
 \widehat{\operatorname{CRPS}}(X,y)=
 \frac{1}{M}\sum_m |X^{(m)}-y|
 -\frac{1}{2M(M-1)}\sum_{m\ne n}|X^{(m)}-X^{(n)}|.
\end{equation}
We additionally report 90\% central interval coverage, width, and interval
score; an energy score for each group's normalized multihorizon trajectory;
and CRPS for the fraction of groups in which at least 20\% of the current cohort
exits. These are evaluations of specified operational projections, not a claim
that a small collection of scores identifies the full high-dimensional law.

To expose decision tradeoffs, let $R_{Ak}=|A|-C_{Ak}$ be the number of current
members that have not exited. For an under/over penalty ratio $r$, select the
$r/(r+1)$ forecast quantile $b$ of $R_{Ak}$ and evaluate
\begin{equation}
 \ell_r(b,R_{Ak})=
 \frac{r\max(R_{Ak}-b,0)+\max(b-R_{Ak},0)}{|A|},
 \qquad r\in\{3,9\}.
\end{equation}
This is a reserve-cost proxy for the current cohort. It is not measured
currency, a staffing policy, or a counterfactual deployment estimate. Future
arrivals and subsequent state returns are outside its scope. Distributional
score improvements and improvements under this particular loss are reported
separately.

We average random-seed repetitions within each forecast origin before paired
time-block resampling. Seeds and individual tasks are not additional independent
temporal observations. The final analysis uses circular moving blocks of one
and two days for the cluster trace and 14 and 28 days for the workflow trace,
reporting when the available span is too short to support a chosen block size.
The resulting intervals are descriptive approximations for finite dependent
traces, rather than guarantees of performance at other sites.
